% ============================================================
% Plantilla: Artículo científico (estructura IMRyD)
% Compila con pdflatex + bibtex.
% ============================================================
\documentclass[11pt,a4paper]{article}

\usepackage[utf8]{inputenc}
\usepackage[T1]{fontenc}
\usepackage{lmodern}
\usepackage[spanish,es-tabla]{babel}
\usepackage[margin=2.5cm]{geometry}
\usepackage{amsmath,amssymb}
\usepackage{graphicx}
\usepackage{booktabs}
\usepackage[hidelinks]{hyperref}

\title{Título del artículo: claro, específico y breve}
\author{
  Primer Autor\thanks{Universidad, correo@ejemplo.com} \and
  Segundo Autor\thanks{Institución, correo@ejemplo.com}
}
\date{\today}

\begin{document}

\maketitle

\begin{abstract}
Resume en un párrafo el objetivo, la metodología, los resultados principales y la
conclusión del trabajo (150 a 250 palabras).

\medskip
\noindent\textbf{Palabras clave:} palabra uno; palabra dos; palabra tres.
\end{abstract}

\section{Introducción}
Presenta el problema, su relevancia y el objetivo del artículo. Cita trabajos previos
\cite{lamport1994}.

\section{Materiales y métodos}
Describe cómo se realizó el estudio para que pueda reproducirse. El error cuadrático
medio se calcula como:
\begin{equation}
  \mathrm{ECM} = \frac{1}{n} \sum_{i=1}^{n} \left( y_i - \hat{y}_i \right)^2
  \label{eq:ecm}
\end{equation}

\section{Resultados}
La Tabla~\ref{tab:comparacion} compara los métodos evaluados según la
ecuación~\eqref{eq:ecm}.

\begin{table}[htbp]
  \centering
  \caption{Comparación de los métodos evaluados.}
  \label{tab:comparacion}
  \begin{tabular}{lrr}
    \toprule
    Método & ECM & Tiempo (ms) \\
    \midrule
    Método A & 0.042 & 120 \\
    Método B & 0.035 & 180 \\
    Propuesto & \textbf{0.028} & 140 \\
    \bottomrule
  \end{tabular}
\end{table}

\section{Discusión}
Interpreta los resultados y compáralos con la literatura \cite{knuth1984}.

\section{Conclusiones}
Resume los aportes del trabajo y propone líneas futuras.

\bibliographystyle{plain}
\bibliography{referencias}

\end{document}
